% !TeX root = ../../Book5.tex
% 纲要标题：### D.1 从微分模到 D-模
% 纲要位置：工作记录/计划纲要.md，第 4908 行。
% * Weyl 代数模、微分算子环及线性方程组
% * 代数微分算子、滤过与伴随分次的符号思想
% * 特征零光滑代数簇上的 D-模与特征簇；明确与线性方程组的联系
% * holonomic D-模的定义和代表例子；holonomic 保留原文，不译作 holomorphic（全纯）；也不是第三卷微分代数的同义词
% #### D.1.1 Weyl 代数模与线性方程组
% #### D.1.2 微分算子、滤过与伴随分次
% #### D.1.3 光滑代数簇上的 D-模与特征簇
% #### D.1.4 holonomic D-模与术语辨析
\section{从微分模到 D-模}
\label{b5:sec:D01}

线性方程 \(P(y)=0\) 中，未知量所在的函数空间可以改变，算子之间的代数关系却可以预先记录下来。D-模把这些关系组成一个模，再通过模同态求解。以下先在仿射直线上计算，随后说明在光滑簇上需要添加哪些几何结构。

\subsection{Weyl 代数模与线性方程组}
\label{b5:subsec:D01:01}

设 \(k\) 为特征零域。在 \(k[x]\) 上，把乘以 \(x\) 的算子仍记为 \(x\)，把求导算子记为 \(\partial\)。Leibniz 法则给出
\[
 \partial x-x\partial=1.
\]
这些算子生成的 \(k\) 代数为一阶 \term[Weyl-daishu]{Weyl 代数}[Weyl algebra] \(A_1(k)\)。它的每个元素唯一写为
\(\sum_{j=0}^m a_j(x)\partial^j\)。存在性来自交换关系不断把 \(x\) 移到左边。若这种表达为零算子，取最大 \(m\)，对该算子反复与乘法算子 \(x\) 取交换子 \(m\) 次，便得乘法算子 \(m!\cdot a_m(x)\) 为零，从而 \(a_m=0\)，归纳得所有系数为零。这也证明抽象生成关系没有额外隐藏关系。

\begin{proposition}\label{b5:prop:D-solution-module}
设 \(k\) 为特征零域，\(D=A_1(k)\)，\(P\in D\)，\(V\) 为左 \(D\) 模。则
\[
 \operatorname{Hom}_D(D/DP,V)\longrightarrow\{v\in V:Pv=0\},
 \qquad f\longmapsto f(1+DP)
\]
是 \(k\) 线性同构。
\end{proposition}
\begin{proof}
左模同态由 \(1+DP\) 的像决定，而 \(P(1+DP)=0\)，故其像必须被 \(P\) 零化。反过来，对满足 \(Pv=0\) 的 \(v\)，令
\(f(Q+DP)=Qv\)。若两个代表元相差 \(RP\)，其作用差为 \(R(Pv)=0\)，所以此式良定义且为左模同态。两个构造互逆。
\end{proof}

对方程组 \(\sum_jP_{ij}y_j=0\)，取自由左模 \(D^r\) 的基 \(e_1,\ldots,e_r\)，再模去由关系 \(\sum_jP_{ij}e_j\) 生成的左子模。同样的论证把该商模到 \(V\) 的同态，识别为 \(V^r\) 中的解。因而系数作用在左边，关系也必须按左模生成，不能随意交换微分算子的次序。

\ProseParagraphBreak
例如 \(D/D\partial\) 的多项式解是常数，\(D/D(\partial-1)\) 的多项式解为零；第二个模本身却非零，在解析函数中还有 \(e^x\) 这样的解。方程模是否为零与某一函数空间里是否有非零解，是不同问题。

\subsection{微分算子、滤过与伴随分次}
\label{b5:subsec:D01:02}

\begin{definition}\label{b5:def:D-operator-order}
设 \(k\) 为特征零域，\(R\) 为交换含幺 \(k\) 代数。在 \(\operatorname{End}_k(R)\) 中递归规定 \(F_{-1}=0\)，并令
\[
 F_m=\bigl\{P:[P,a]\in F_{m-1}\text{ 对所有 }a\in R\bigr\}\qquad(m\ge0),
\]
其中 \(P\in\operatorname{End}_k(R)\)，\(a\) 表示乘以 \(a\) 的算子，\([P,a]=Pa-aP\)。\(\mathcal D(R)=\bigcup_{m\ge0}F_m\) 的元素称为 \(R\) 上的
\term[daishu-weifen-suanzi]{代数微分算子}[algebraic differential operator]，最小的相应 \(m\) 称为其阶。
\end{definition}

\(F_0=\operatorname{End}_R(R)=R\)，所以零阶算子就是乘法。恒等式
\([PQ,a]=P[Q,a]+[P,a]Q\) 归纳给出 \(F_mF_n\subset F_{m+n}\)。在 \(R=k[x]\) 中，归纳还给出
\[
 F_m=\sum_{j=0}^m k[x]\partial^j.
\]
事实上，若 \([P,x]=\sum_{j<m}b_j(x)\partial^j\)，令
\(Q=\sum_{j<m}b_j(x)\partial^{j+1}/(j+1)\)，则
\([P-Q,x]=0\)。因此 \(P-Q\) 与一切多项式乘法交换，必为乘法算子。这证明 \(\mathcal D\bigl(k[x]\bigr)=A_1(k)\)。

\ProseParagraphBreak
这个阶滤过把 \(x\) 放在零阶，把 \(\partial\) 放在一阶。其伴随分次环为
\[
 \operatorname{gr}_F D\simeq k[x,\xi],
\]
其中 \(\xi\) 是 \(\partial\) 的一阶符号；对
\(P=a_m(x)\partial^m+\cdots\)、\(a_m\ne0\)，主符号是 \(a_m(x)\xi^m\)。因为交换子降低阶，原来的非交换关系在伴随分次环中变为交换关系。

\ProseParagraphBreak
还有常用于增长估计的 Bernstein 滤过，把 \(x,\partial\) 都放在一次。它同样产生多项式分次环，但滤过不相同。例如 \(\partial-x\) 的阶主符号为 \(\xi\)，Bernstein 主符号却为 \(\xi-x\)。以下几何特征簇始终用阶滤过；不能因两个分次环抽象同构就把两个主符号混用。两种滤过的区别亦见 \cite[Section 3.2, pp. 88--90]{HottaTakeuchiTanisaki2008DModules}。

\subsection{光滑代数簇上的 D-模与特征簇}
\label{b5:subsec:D01:03}

设 \(X\) 为光滑复代数簇，\(\mathcal O_X\) 为正则函数层，\(\mathcal T_X=\operatorname{Der}_{\mathbb C}(\mathcal O_X)\) 为切向量层。层 \(\mathcal End_{\mathbb C}(\mathcal O_X)\) 的局部截面是与开集限制相容的 \(\mathbb C\) 线性层态射。把函数视为乘法算子，由 \(\mathcal O_X\) 与 \(\mathcal T_X\) 在其中生成的子代数层记为 \(\mathcal D_X\)，称为微分算子层。这一定义直接在层态射中进行，因而其算子已经具有限制映射。

\ProseParagraphBreak
局部可以选择 étale 坐标 \(x_1,\ldots,x_n\)，使 \(\mathrm{d}x_1,\ldots,\mathrm{d}x_n\) 为微分层的基。其对偶导子 \(\partial_1,\ldots,\partial_n\) 满足
\[
 [\partial_i,a]=\frac{\partial a}{\partial x_i},\qquad
 [\partial_i,\partial_j]=0.
\]
把函数移到左边、把相互交换的导子按次序排列，便得到局部表达
\[
 \mathcal D_U=\bigoplus_{\alpha\in\mathbb N^n}\mathcal O_U\partial^\alpha,
 \qquad
 F_m\mathcal D_U=\bigoplus_{|\alpha|\le m}\mathcal O_U\partial^\alpha.
\]
其中直和表达的唯一性、阶滤过与递归交换子定义的一致性，以及
\(\operatorname{gr}\mathcal D_X\simeq\operatorname{Sym}_{\mathcal O_X}\mathcal T_X\)，可由以下符号计算得到。右侧是余切丛 \(T^*X\) 上沿纤维的多项式函数。
这个对偶方向来自求值：切向量场 \(v\) 给出余切丛上的函数 \((x,\xi)\mapsto\xi(v_x)\)，它在每个余切纤维上为线性函数。若 \(\xi=\sum_i\xi_i\,\mathrm{d}x_i\)，则 \(\xi(\partial_i)=\xi_i\)；因此导子 \(\partial_i\) 的符号恰是余切坐标 \(\xi_i\)，其对称乘积给出纤维上的多项式函数。
\begin{proofsketch}
对于由函数和导子组成的表达，反复与各个 \(x_i\) 取交换子：若最高总次数为 \(m\)，任取 \(|\alpha|=m\)，与 \(x_i\) 分别取 \(\alpha_i\) 次交换子所得的算子就是乘以 \(\alpha!a_\alpha\)。因此零算子的最高系数全为零，归纳证明正规表达唯一。交换子 \([P,a]\) 降低至少一阶，故上面的阶滤过满足\cref{b5:def:D-operator-order}的递归条件。

反过来，设层态射 \(P\) 按递归交换子条件的阶数至多为 \(m\)。其 \(m\) 重交换子是乘法算子；所乘函数关于这 \(m\) 个函数参数是对称的多导子。这由乘法算子彼此交换及 Leibniz 恒等式得到。由于 \(\mathrm{d}x_i\) 为微分层的基，该多导子由在坐标参数上的取值决定。用这些取值除以相应的 \(\alpha!\)，得到系数 \(a_\alpha\)，使
\(P-\sum_{|\alpha|=m}a_\alpha\partial^\alpha\) 的全部 \(m\) 重交换子为零，即阶数至多为 \(m-1\)。归纳后 \(P\) 属于函数与导子生成的算子层，证明两种定义一致。

令 \(\xi_i\) 为 \(\partial_i\) 的一阶符号，局部便有 \(\operatorname{gr}\mathcal D_U=\mathcal O_U[\xi_1,\ldots,\xi_n]\)。换坐标时，一阶符号按切向量变换，各阶乘积按其对称幂变换，故这些识别粘合成所示对称代数同构。局部坐标表达、阶滤过性质、递归描述和主符号识别分别见\cite[Section 1.1, Proposition 1.1.3, Exercise 1.1.4 and the principal-symbol discussion, pp. 15--17]{HottaTakeuchiTanisaki2008DModules}。
\end{proofsketch}

求出某个函数空间中的全部解，往往比记录方程的最高阶约束困难。阶滤过使最高阶关系落在交换的符号代数中，因而可以用其支撑描述约束所在的余切方向。对一个模需要同时记录全部关系，并保证所选滤过由有限数据控制；良滤过提供了这个条件。

\ProseParagraphBreak
在仿射坐标开集上，\(\operatorname{gr}\mathcal D_X\) 的截面环为 Noether 多项式环。左理想或右理想的分次理想都有有限个齐次生成元，提升这些生成元并按阶归纳，便得到原理想的有限生成元。因此局部算子环及各茎均左右 Noether。对于底层 \(\mathcal O_X\) 模拟凝聚、且在 \(\mathcal D_X\) 上局部有限生成的模，局部展示的核仍在 \(\mathcal O_X\) 上拟凝聚；在仿射开集取截面后，Noether 性给出核的有限生成元。因而 \(\mathcal D_X\) 是相干环层，这类模也局部有限展示。反过来，局部有限展示的 \(\mathcal D_X\) 模是有限自由 \(\mathcal D_X\) 模之间态射的余核，所以其底层 \(\mathcal O_X\) 模拟凝聚。这说明下面的有限展示定义与通常的相干 D-模概念一致；参见\cite[Propositions 1.4.6 and 1.4.9, pp. 26--27]{HottaTakeuchiTanisaki2008DModules}。

\begin{definition}\label{b5:def:D-good-characteristic}
设 \(X\) 为光滑复代数簇，\(\mathcal O_X\) 为正则函数层，\(\mathcal D_X\) 为 \(\mathcal O_X\) 上的代数微分算子层，\(F_i\mathcal D_X\) 为阶数不超过 \(i\) 的算子层（\(i<0\) 时取零）。\(\mathcal D_X\) 的左模层称为左 D-模；若它局部有限展示，则称为相干 D-模。设 \(M\) 为相干左 D-模，\((F_jM)_{j\in\mathbb Z}\) 为递增、穷尽、局部下有界的滤过。若各 \(F_jM\) 为相干 \(\mathcal O_X\) 模，
\(F_i\mathcal D_XF_jM\subset F_{i+j}M\)，且
\(\operatorname{gr}_F M\) 在 \(\operatorname{gr}\mathcal D_X\) 上局部有限生成，则称为
\term[liang-lvguo]{良滤过}[good filtration]。
通过 \(\operatorname{gr}\mathcal D_X\simeq\operatorname{Sym}_{\mathcal O_X}\mathcal T_X\)，其中 \(\mathcal T_X\) 为切向量层，将 \(\operatorname{gr}_F M\) 视为余切丛 \(T^*X\) 上的相干层。其支撑集称为 \(M\) 的
\term[tezheng-cu]{特征簇}[characteristic variety]，记为 \(\operatorname{Ch}(M)\)。
\end{definition}

\begin{proposition}[良滤过与特征簇]\label{b5:prop:D-good-filtration-invariance}
设 \(X\) 为光滑复代数簇，\(\mathcal D_X\) 为代数微分算子层，\(M\) 为相干左 \(\mathcal D_X\) 模。它局部具有关于阶滤过的良滤过，以其伴随分次模在 \(T^*X\) 上的支撑定义的特征簇 \(\operatorname{Ch}(M)\) 不依赖良滤过的选择。对相干左 \(\mathcal D_X\) 模的短正合列
\(0\to M'\to M\to M''\to0\)，有
\(\operatorname{Ch}(M)=\operatorname{Ch}(M')\cup\operatorname{Ch}(M'')\)。若 \(\pi:T^*X\to X\) 为丛投影，则还有
\[
 \operatorname{Supp}_{\mathcal O_X}(M)=\pi\bigl(\operatorname{Ch}(M)\bigr),
\]
且这个支撑是闭集。
\end{proposition}
\begin{proof}
局部取有限个 \(\mathcal D_X\) 生成元 \(u_i\)，指定它们的整数滤过次数 \(d_i\)，令
\(F_jM=\sum_iF_{j-d_i}\mathcal D_Xu_i\)。每项在 \(\mathcal O_X\) 上相干，分次模由这些元的符号生成，故为良滤过。

两个良滤过 \(F,G\) 的有限个分次生成元提升到模中，分别在另一滤过的有限次数出现。因此可取 \(c\ge0\)，使
\(F_jM\subseteq G_{j+c}M\)、\(G_jM\subseteq F_{j+c}M\) 对所有 \(j\) 成立。设次数 \(m\) 的齐次符号 \(a\) 零化 \(\operatorname{gr}_FM\)，取它的算子提升 \(P\)，则
\(PF_jM\subseteq F_{j+m-1}M\)。对 \(N>2c\)，得到
\[
 P^NG_jM\subseteq F_{j+c+Nm-N}M
          \subseteq G_{j+Nm+2c-N}M
          \subseteq G_{j+Nm-1}M.
\]
所以 \(a^N\) 零化 \(\operatorname{gr}_GM\)。交换 \(F,G\) 后得到反向包含，故两分次模的零化理想有相同根理想，支撑相同。

在 \(M'\) 上取 \(M\) 的交滤过，在 \(M''\) 上取商滤过，得到分次短正合列。分次环局部为 Noether 多项式环，所以分次子模和商模都有限生成，各滤过项也相干；这些仍为良滤过。交换环上短正合列的支撑是两端支撑的并，故得到结论。参见\cite[Theorems 2.1.3 and 2.2.1, pp. 58--59]{HottaTakeuchiTanisaki2008DModules}。

最后证明支撑与投影的关系。局部记 \(G=\operatorname{gr}_FM\)。滤过穷尽且下有界，故 \(M_x=0\) 当且仅当 \(G_x=0\)。\(G_x\) 在 \(\mathcal O_{X,x}[\xi_1,\ldots,\xi_n]\) 上分次有限生成，各次分量在 \(\mathcal O_{X,x}\) 上有限生成；分次 Nakayama 论证于是给出
\[
 G_x\ne0\quad\Longleftrightarrow\quad
 G_x/(\mathfrak m_x,\xi_1,\ldots,\xi_n)G_x\ne0.
\]
具体说，若右侧为零，最低非零次数的分量便等于其极大理想倍，由普通 Nakayama 引理只能为零，矛盾。右侧非零等价于 \((x,0)\in\operatorname{Ch}(M)\)。特征簇由齐次理想定义，沿纤维伸缩不变；闭性使 \((x,\xi)\) 在其中时 \((x,0)\) 也在其中。因此 \(\pi(\operatorname{Ch}(M))=\{x:(x,0)\in\operatorname{Ch}(M)\}\)，并与 \(\operatorname{Supp}_{\mathcal O_X}(M)\) 一致；这个集合是闭集。
\end{proof}

这里相干性是对 \(\mathcal D_X\) 而言，不要求 \(M\) 在 \(\mathcal O_X\) 上有限生成。

\begin{proposition}\label{b5:prop:D-principal-characteristic}
设 \(X=\mathbb A^1_{\mathbb C}\)，\(D=A_1(\mathbb C)\) 取阶滤过，\(\xi\) 为 \(\partial\) 的一阶符号。设非零算子 \(P=a_m(x)\partial^m+\cdots\in D\)，其中 \(m\ge0\)、\(a_m\in\mathbb C[x]\setminus\{0\}\)。对左商模 \(M=D/DP\) 的商滤过，有
\[
 \operatorname{gr}M\simeq\mathbb C[x,\xi]/\bigl(a_m(x)\xi^m\bigr),
 \qquad
 \operatorname{Ch}(M)=V\bigl(a_m(x)\xi^m\bigr).
\]
\end{proposition}
\begin{proof}
主符号的乘积满足 \(\sigma(QP)=\sigma(Q)\sigma(P)\)：最高阶系数的乘积非零，故阶相加。左理想 \(DP\) 的每个元素就是某个 \(QP\)，所以其分次理想恰由 \(\sigma(P)=a_m(x)\xi^m\) 生成。商滤过的分次模于是为所写商环，它的支撑为该理想的零点集。
\end{proof}

例如 \(D/D\partial\) 的特征簇是零截面 \(\xi=0\)，\(D/Dx\) 的特征簇是点 \(x=0\) 上的整条余切纤维。后一模非零：在 \(\mathbb C[z]\) 上令 \(\partial\) 乘以 \(z\)、\(x=-\mathrm{d}/\mathrm{d}z\)，可验证 \(\partial x-x\partial=1\)，且向量 \(1\) 被 \(x\) 零化。这给出了该商模的具体模型。

\ProseParagraphBreak
再取 \(P=x\partial-\lambda\)、\(\lambda\in\mathbb C\)。它的阶主符号为 \(x\xi\)，所以
\[
\operatorname{Ch}\bigl(D/D(x\partial-\lambda)\bigr)
=V(x\xi)=V(\xi)\cup V(x).
\]
在 \(x\ne0\) 处只留下零截面；在原点，最高阶系数消失，整条余切纤维出现。不同的 \(\lambda\) 都给出同一特征簇，而方程 \(xy'=\lambda y\) 在避开原点的单连通解析区域内有解 \(y=cx^\lambda\)，其绕圈因子随 \(\lambda\) 改变。特征簇保留了最高阶约束的几何形状，低阶参数及解析延拓还需要其他资料。

\subsection{holonomic D-模与术语辨析}
\label{b5:subsec:D01:04}

特征簇越小，方程的约束通常越强。不过约束不能任意增加而仍保留非零相干模。

\begin{theorem}[Bernstein 不等式]\label{b5:thm:D-bernstein}
设 \(n\ge0\)，\(X\) 为纯 \(n\) 维光滑复代数簇，\(\mathcal D_X\) 为代数微分算子层，\(M\ne0\) 为相干左 \(\mathcal D_X\) 模，\(\operatorname{Ch}(M)\subseteq T^*X\) 为其关于阶滤过的特征簇。则
\(\dim\operatorname{Ch}(M)\ge n\)。
\end{theorem}

\begin{proofsketch}
按 \(n=\dim X\) 归纳，\(n=0\) 时结论直接成立。光滑簇的不可约分量彼此不交，可以先限制到 \(M\) 非零的一个分量并仍记它为 \(X\)。由\cref{b5:prop:D-good-filtration-invariance}，模的支撑等于特征簇的投影。若支撑为 \(X\)，特征簇便包含零截面，维数至少为 \(n\)。否则，支撑为非空真闭子集。取其中一个不可约分量的光滑点附近的开集，并去掉其他支撑分量，便可选择光滑超曲面 \(S=\{x_1=0\}\) 包含所剩支撑；限制后的 \(M\) 仍非零。取与 \(S\) 相容的 étale 坐标及法向算子 \(\partial_1\)。底层 \(\mathcal O_X\) 模拟凝聚，所以支撑在 \(S\) 上意味着局部化 \(M[x_1^{-1}]=0\)，每个局部元素 \(m\) 都被某个 \(x_1\) 的幂零化。

令 \(N=\ker(x_1:M\to M)\)，把它视为 \(S\) 上的层。函数在 \(N\) 上通过 \(\mathcal O_S=\mathcal O_X/(x_1)\) 作用，切向导子 \(\partial_2,\ldots,\partial_n\) 保持这个核，因而赋予它左 \(\mathcal D_S\) 模结构。对非零局部元素 \(m\)，取最小的 \(r\ge1\) 使 \(x_1^rm=0\)；则 \(x_1^{r-1}m\) 是 \(N\) 的非零元素。Weyl 关系给出逐元素有限求和的投影
\[
 p(m)=\sum_{j\ge0}\frac{\partial_1^j x_1^jm}{j!}\in N,
 \qquad
 m=\sum_{j\ge0}\frac{(-1)^j}{j!}\partial_1^j p(x_1^jm).
\]
这个投影是在取法向导数变量的常数项。确实，对 \(n\in N=\ker x_1\)，Weyl 关系递推给出
\(x_1^j\partial_1^r n=(-1)^j r!\,\partial_1^{r-j}n/(r-j)!\)（\(0\le j\le r\)），\(j>r\) 时为零。所以
\[
p(\partial_1^r n)
=\left(\sum_{j=0}^r(-1)^j\binom rj\right)\partial_1^r n
=\begin{cases}n,&r=0,\\0,&r>0.\end{cases}
\]
它保留 \(N\) 中的常数项，并消去全部正次法向导数。对一般 \(m\)，第一式确实被 \(x_1\) 零化，因为相邻项用
\(x_1\partial_1^j=\partial_1^jx_1-j\partial_1^{j-1}\) 相消；第二式由代入第一式并合并有限和得到。若 \(\sum_{j=0}^s\partial_1^jn_j=0\)、\(n_j\in N\)，作用 \(x_1^s\) 得 \((-1)^ss!n_s=0\)，再按 \(s\) 下降归纳。因此这些和唯一，局部有
\[
 M\simeq\mathbb C[\partial_1]\otimes_{\mathbb C}i_*N,
 \qquad i:S\hookrightarrow X,
\]
其中 \(x_1\) 按 \(-\mathrm{d}/\mathrm{d}\partial_1\) 作用。这是 Kashiwara 等价在超曲面支撑情形的局部模型，参见\cite[Theorem 1.6.1 and its proof, pp. 48--50]{HottaTakeuchiTanisaki2008DModules}。

为核对相干性，取有限个局部 \(\mathcal D_X\) 生成元 \(m_1,\ldots,m_q\)，并选择 \(r_i\ge1\) 使 \(x_1^{r_i}m_i=0\)。令 \(N_0\) 为有限个元素
\(p(x_1^jm_i)\)、\(1\le i\le q\)、\(0\le j<r_i\)，在 \(\mathcal D_S\) 上生成的子模，并置
\(W=\sum_{j\ge0}\partial_1^ji_*N_0\)。Taylor 式使所有 \(m_i\) 属于 \(W\)。这个子模对 \(x_1\)、\(\partial_1\) 及切向导子封闭；对任意局部函数 \(a\)，恒等式
\[
 a\partial_1^jn
 =\sum_{\ell=0}^j(-1)^\ell\binom{j}{\ell}
   \partial_1^{j-\ell}\bigl((\partial_1^\ell a)n\bigr)
 \qquad(n\in N_0)
\]
也证明它对函数作用封闭，其中括号内的函数在 \(S\) 上限制后作用于 \(N_0\)。因此 \(W=M\)。上述唯一分解又给出 \(W\cap\ker x_1=i_*N_0\)，故 \(N=N_0\) 在 \(\mathcal D_S\) 上局部有限生成。其底层 \(\mathcal O_S\) 模拟凝聚，因为它是 \(\mathcal O_X\) 拟凝聚模之间映射的核且被 \(x_1\) 零化；由前面的 Noether 性与相干性结论，\(N\) 是相干左 \(\mathcal D_S\) 模。

给 \(N\) 取良滤过 \(G\)，在 \(M\) 上令
\[
 F_jM=\sum_{r\ge0}\partial_1^ri_*G_{j-r}N.
\]
每个固定 \(j\) 只有有限个非零项。上面的函数作用公式与法向关系说明这是良滤过，并有
\[
 \operatorname{gr}_FM\simeq
 \mathbb C[\xi_1]\otimes_{\mathbb C}i_*\operatorname{gr}_GN.
\]
这里 \(x_1\) 的符号作用为零，\(\xi_1\) 是自由的法向余切坐标；切向符号按 \(\operatorname{gr}_GN\) 的原作用给出。因此
\(\dim\operatorname{Ch}(M)=\dim\operatorname{Ch}(N)+1\)。归纳假设给出 \(\dim\operatorname{Ch}(N)\ge n-1\)，于是 \(\dim\operatorname{Ch}(M)\ge(n-1)+1=n\)，归纳完成。局部模型在重叠坐标上的层论识别由 Kashiwara 等价给出；闭嵌入的特征簇公式及上述支撑归纳见\cite[Lemma 2.3.5 and the proof of Corollary 2.3.2, pp. 63--64]{HottaTakeuchiTanisaki2008DModules}。
\end{proofsketch}

设 \(X\) 为纯 \(n\) 维光滑复代数簇，\(\mathcal D_X\) 为代数微分算子层。零左 \(\mathcal D_X\) 模，以及满足 \(\dim\operatorname{Ch}(M)=n=\dim X\) 的非零相干左 \(\mathcal D_X\) 模，称为
\term[holonomic-D-mo]{holonomic D-模}[holonomic D-module]。
因此\cref{b5:prop:D-principal-characteristic}中，凡商模非零且 \(P\ne0\)，其特征簇在平面里是一条曲线，所得模便是 holonomic；反之自由模 \(D\) 的特征簇是整个二维余切丛，故不是 holonomic。

\ProseParagraphBreak
有限秩向量丛上的平坦联络给出另一类例子。在局部基中，把 \(\partial_i\) 作用写成 \(\nabla_i=\partial_i+B_i\)，Leibniz 法则处理它与函数的关系，而
\[
 [\nabla_i,\nabla_j]=0
 \quad\Longleftrightarrow\quad
 \partial_iB_j-\partial_jB_i+[B_i,B_j]=0
\]
保证各偏导作用满足 D-模关系。这就是平坦性。取 \(F_jM=M\) 对 \(j\ge0\)、\(F_jM=0\) 对 \(j<0\)，便得良滤过，其正次符号全部作用为零，特征簇为零截面，所以非零联络是 holonomic。一般 holonomic 模却未必是向量丛上的联络，\(D/Dx\) 就说明了点支撑的可能性。

\ProseParagraphBreak
英文 holonomic 表示这里的特征簇维数条件，holomorphic 则表示复解析意义的“全纯”，二者不能互译。D-模与第三卷的微分代数也有不同的研究对象：前者研究线性微分算子作用及其模，后者还允许非线性微分多项式关系；两者的联系须在具体对象上讨论。
